\documentclass{article}

% Paquetes para soporte de español y márgenes
\usepackage[utf8]{inputenc}
\usepackage[spanish]{babel}
\usepackage{geometry}

% Configuración de la página
\geometry{a4paper, margin=1in}

\title{Cuestionario Semana 6: El Diodo: Fundamentos y Rectificación}
\author{Plan de Estudios de Electrónica}
\date{\today}

\begin{document}

\maketitle

\section*{Preguntas de Repaso}

\begin{enumerate}
    \item ¿Cuáles son las diferencias fundamentales en el comportamiento y las características i-v (corriente-voltaje) entre un diodo ideal y un diodo real?
    
    \item Explica el funcionamiento de un circuito rectificador de media onda. Dibuja el esquema del circuito y la forma de onda de salida si la entrada es una señal sinusoidal.
    
    \item Describe el modelo exponencial del diodo. ¿Qué representa la ecuación de Shockley y por qué es importante para entender el comportamiento del diodo en polarización directa?
    
    \item ¿Qué ventajas ofrece un rectificador de onda completa (ya sea con transformador de toma intermedia o en puente) en comparación con un rectificador de media onda?
    
    \item ¿En qué consiste el modelo de "tramos lineales" (piecewise linear model) para un diodo y en qué situaciones de análisis de circuitos resulta más útil que el modelo ideal o el exponencial?
\end{enumerate}

\end{document}